% revised_text.tex
% Drop-in LaTeX for the revised sections. Citation keys match references.bib.
% Uses \citep / \citet (natbib, author-year) to match the draft's APA-style
% in-text citations. Swap to \cite if the venue template is numeric.
% Needs: \usepackage{natbib,amsmath,booktabs,graphicx}
%
% Items in [SQUARE BRACKETS] are facts only the authors know. Fill them in.

% ===========================================================================
% I. INTRODUCTION, last paragraph (comments 2-4)
% ===========================================================================

This paper presents a preliminary approach to assessing handwritten scientific
documents for machine readability before OCR transcription. The section that
follows provides a brief background on the use of analog notebooks in labs,
defines Computational Archival Science (CAS), and details how it relates to our
research. Next, we state our research goals and describe the methods used to
compute a readability score from three measurements (baseline consistency,
spatial organization, and text fragmentation). As part of our approach, the
score is then compared to the character error rate (CER) of an OCR
transcription, to check whether the score tracks actual transcription error.
Finally, we discuss research implications and future work, identify next
steps, and present our conclusions.

% ===========================================================================
% II. BACKGROUND, citation fixes (comment 6)
% ===========================================================================
% "...(Marciano et al., 2018)"  ->  \citep{marciano2018archival}
% "...(Jenkinson, 1922)"         ->  \citep{jenkinson1922manual}
% Both were missing from the reference list and are now in references.bib.

% ===========================================================================
% IV.A DATA (comments 8-11, 38, 49, plus ground truth and OCR confidence)
% ===========================================================================

\subsection{Data}

This preliminary study uses three pages: one representative page
($1455 \times 1980$ pixels) from each of three student lab notebooks. A MOF
researcher visually assessed the scans and grouped them by handwriting style
as neat, average, or messy (Fig.~1).

Each page was prepared (Section~IV-B) and three structural measurements were
taken from the prepared image: baseline consistency ($b$), spatial
organization ($s$), and text fragmentation ($f$). These were combined into a
readability score from 0 (low readability) to 1 (high readability). To check
whether the score reflects real transcription difficulty, each original,
unprepared scan was also transcribed with Tesseract OCR
\citep{smith2007tesseract, tesseract_manual} through pytesseract
\citep{pytesseract}, and the output was scored against a ground-truth
transcription.

% --- Ground truth (checklist item 1; comment 49) ---
The ground truth for each page is a manual transcription typed by
[NAME / ROLE, e.g., ``the first author'' or ``a MOF researcher''], who read
each scan and typed its text into a plain UTF-8 text file, keeping the page's
line breaks and reading order and entering table rows as space-separated
values [ADJUST IF NEEDED: e.g., ``chemical structure drawings were recorded by
their labels only'']. No OCR output was used to produce the ground truth. The
resulting references contain 997 (neat), 463 (average), and 845 (messy)
characters.

% --- CER and JiWER (checklist item 1; comment 38) ---
CER was computed with the JiWER library \citep{jiwer} as the character-level
edit distance (substitutions, deletions, and insertions) between the OCR
output and the ground truth, divided by the number of characters in the
ground truth. Case and punctuation were kept, and no other text normalization
was applied.
% OPTIONAL, see revision_notes.md section A.1: "Line breaks in the ground truth
% were compared against the single spaces that join OCR words, which raises
% CER slightly; collapsing all whitespace lowers CER by 0.2 to 2.8 points and
% does not change the page ranking."

% --- OCR confidence (checklist item 5; comments 39-40) ---
OCR confidence is the mean of the word-level confidence values (0 to 100)
that Tesseract reports through \texttt{pytesseract.image\_to\_data}, averaged
over every word on the page and excluding the $-1$ placeholders Tesseract
assigns to non-word layout entries. Tesseract [VERSION, e.g., 5.x] was run
with its default settings (English model, automatic page segmentation).

% ===========================================================================
% IV.B IMAGE PREPARATION (comments 12-14 citations)
% ===========================================================================
% First paragraph: keep as is. Remove "then processed" from the Data section
% (done above) so "preparation" is used consistently (comment 12).

The first preparation stage, grayscale conversion, reduces each pixel from
three color channels to a single intensity value, which simplifies the
analysis when ink color is not being evaluated \citep{gonzalez2018dip}. A
$3 \times 3$ Gaussian blur then smooths small pixel-level variations
\citep{gonzalez2018dip, opencv_filtering}. Adaptive thresholding separates the
handwriting from the background by computing a separate threshold for each
pixel from a Gaussian-weighted $31 \times 31$ neighborhood (offset $C = 11$)
\citep{opencv_misc_transforms}. Unlike global thresholding, which applies one
threshold to the whole image, adaptive thresholding follows local changes in
illumination, which makes it well suited to scanned documents
\citep{sauvola2000binarization, bradley2007adaptive, binmakhashen2019layout}.
The threshold is inverted, so ink becomes white (foreground) and paper becomes
black (background); OpenCV's morphology and connected-component functions
treat nonzero (white) pixels as the object being measured
\citep{opencv_morph_tutorial, opencv_shape}.

Morphological processing controls the shape and connectivity of the
foreground regions \citep{serra1982morphology, haralick1987morphology}.
Morphological opening (erosion followed by dilation, $3 \times 3$ square
kernel, one pass) removes isolated foreground specks smaller than the kernel,
which lowers the chance that image noise is counted as handwriting
\citep{haralick1987morphology, gonzalez2018dip}. Morphological closing, the
reverse operation (dilation followed by erosion), joins nearby foreground
strokes \citep{haralick1987morphology}. Closing ($5 \times 5$ kernel, two
passes) is applied only when measuring text fragmentation (Section~IV-C).

% --- Fig. 2 caption (it shows script output, not image preparation) ---
% \caption{Readability scoring output (Python 3.13) for the neat notebook
% page: baseline angle deviation ($\sigma_\theta$), the three metric scores,
% the combined score, and the Table~I class of each metric.}

% --- Fig. 3 caption (checklist item 6; comments 15-20) ---
% RECOMMENDED: rearrange the five panels in processing order, left to right,
% top row then bottom row, and label them (a)-(e). Then use:
\begin{figure}[t]
  \centering
  % \includegraphics[width=\linewidth]{fig3_preparation.png}
  \caption{Image preparation for the neat notebook page.
  (a) Grayscale conversion of the original scan.
  (b) Adaptive thresholding (inverted) followed by morphological opening; ink
      is white because OpenCV measures nonzero pixels as foreground.
  (c) Canny edge map of (b), shown for illustration only; no metric uses it.
  (d) Morphological closing of (b), which joins nearby strokes before text
      fragmentation is measured.
  (e) Connected components of (d), each shown in a different color.}
  \label{fig:preparation}
\end{figure}
% If the panels are NOT rearranged, the current layout reads:
% top row, left to right: (b), (d), (e); bottom row, left to right: (c), (a).
% All arrows in the current figure point right, so "from right to left" is
% not correct for either row.

% ===========================================================================
% IV.C READABILITY MEASUREMENTS (comments 21-33)
% ===========================================================================

\subsection{Readability Measurements}

Handwritten documents vary in baseline position, inter-line spacing,
character size, and character overlap, all of which complicate text-line
segmentation and recognition
\citep{papavassiliou2010segmentation, rabaev2026baseline}. After image
preparation, three structural metrics are computed: baseline consistency
($b$), spatial organization ($s$), and text fragmentation ($f$). Each raw
measurement is converted to a 0--1 score by comparing it with a fixed
reference value: the value at which the score reaches 0 (or at which a term
reaches its maximum). Scores outside 0--1 are clipped. Higher scores always
mean higher readability. The reference values were set by the authors'
judgement, not derived from data, and are listed in Table~\ref{tab:reference}.

\begin{itemize}
  \item \textbf{Baseline consistency ($b$).} A horizontal projection profile
  (the count of ink pixels in each row) locates text bands
  \citep{likforman2007textline, papavassiliou2010segmentation}. Connected
  components (minimum 15 pixels) are assigned to the band containing their
  center, and bands with fewer than four components are skipped. For each
  band, a least-squares line is fitted through the bottom-center points of the
  components \citep{ogorman1993docstrum, opencv_shape}, and fits with a mean
  residual above 8 pixels are rejected. The baseline angle deviation
  $\sigma_\theta$ is the standard deviation of these line angles, weighted by
  the number of components in each line, and
  \[ b = 1 - \frac{\sigma_\theta}{15^\circ}. \]
  A page whose lines all share one angle scores 1; a spread of $15^\circ$ or
  more scores 0.

  \item \textbf{Spatial organization ($s$).} Rows whose ink count exceeds 15\%
  of the page's maximum row are grouped into bands, and the vertical distance
  between consecutive band centers is measured
  \citep{likforman2007textline}. With $\mathrm{Var}(d)$ the variance of these
  spacings and $H$ the page height in pixels,
  \[ s = 1 - \frac{\mathrm{Var}(d)}{H^2/16}. \]
  The reference $H^2/16$ corresponds to a spacing standard deviation of one
  quarter of the page height (495 pixels for these pages).

  \item \textbf{Text fragmentation ($f$).} After morphological closing,
  connected components smaller than 15 pixels are discarded as noise. Three
  terms are computed, each oriented so that a larger value means more
  fragmentation, and each limited to 0--1:
  \begin{itemize}
    \item component density $D = \min(40 \cdot n / A, 1)$, where $n$ is the
    number of components and $A$ is their total ink area in pixels
    \citep{fujisawa1992segmentation};
    \item nearest-neighbor dispersion $N = \min(\bar{r} / \sqrt{HW}, 1)$,
    where $\bar{r}$ is the mean distance from each component's centroid to
    its nearest neighbor and $\sqrt{HW}$ is the square root of the page area
    \citep{ogorman1993docstrum, clark1954nearest};
    \item small-component ratio $R$, the share of components under 150 pixels,
    such as broken strokes, isolated punctuation, and residual noise
    \citep{fujisawa1992segmentation, ye2013dqa}.
  \end{itemize}
  These are combined into a fragmentation index and inverted, so that higher
  $f$ means less fragmentation:
  \[ f = 1 - (0.45\,D + 0.35\,N + 0.20\,R). \]
  The weights (0.45, 0.35, 0.20) reflect the authors' judgement of each
  term's likely effect on OCR; we found no published weights for these
  terms, and this is treated as a limitation (Section~VI).
\end{itemize}

\begin{table}[t]
  \centering
  \caption{Reference values used to scale each measurement to 0--1. All were
  set by judgement.}
  \label{tab:reference}
  \begin{tabular}{@{}lll@{}}
    \toprule
    Metric / term & Raw measurement & Reference value (score or term limit) \\
    \midrule
    $b$ & $\sigma_\theta$ (degrees) & $15^\circ$ ($b = 0$ at or above) \\
    $s$ & Var. of band spacing (px$^2$) & $H^2/16 = 245{,}025$ px$^2$ (SD 495 px) \\
    $f$: $D$ & components per ink pixel & $1/40$ ($D = 1$ at or above) \\
    $f$: $N$ & mean nearest-neighbor distance (px) & $\sqrt{HW} = 1697$ px \\
    $f$: $R$ & share of components $<150$ px & none (already a proportion) \\
    \bottomrule
  \end{tabular}
\end{table}

% ===========================================================================
% IV.D COMBINED READABILITY SCORE (comments 30-35)
% ===========================================================================

\subsection{Combined Readability Score}

The combined score is a weighted average of the three metrics,
\[ \text{Score} = 0.40\,b + 0.20\,s + 0.40\,f. \]
Because the weights sum to 1 and each metric lies in 0--1, the score also lies
in 0--1. The weights reflect the authors' judgement that baseline consistency
and fragmentation affect OCR more than line spacing does; no prior work
supplies weights for these metrics. Each metric and the combined score are
labelled High ($\geq 0.80$), Moderate (0.50--0.79), or Low ($< 0.50$)
(Table~I). These cut-offs were also chosen by judgement.

% Table I caption:
% \caption{Classification bands applied to each metric and to the combined
% score.}
% Also fix the middle row of Table I: "0.50-0.75" -> "0.50-0.79".

Finally, the readability score was compared with the CER and with OCR
confidence using Pearson's $r$ and Spearman's $\rho$.

% ===========================================================================
% V. RESULTS  (make this its own section; comments 36-48)
% ===========================================================================

\section{Results}

Table~\ref{tab:raw} lists every measurement for the three pages. All combined
scores fell in the High band, between 0.920 and 0.924. The pages ranked in the
same order by combined score as by CER (neat, average, messy), giving Pearson
$r = -0.882$ and Spearman $\rho = -1.000$ (Table~II). OCR confidence was
positively correlated with CER ($r = 0.674$; $\rho = 0.500$), and the
correlation between the readability score and OCR confidence was
$r = -0.247$, $\rho = -0.500$.

With three pages, none of these correlations is statistically significant
(combined score vs.\ CER: Pearson $p = 0.31$; Spearman exact two-sided
$p = 0.33$). With $n = 3$ the smallest attainable Spearman $p$-value is 0.33,
so these correlations describe the ranking of these three pages only and are
not evidence of a general relationship.

\begin{table}[t]
  \centering
  \caption{Measurements for each page. $\sigma_\theta$ and spacing SD are the
  raw values behind $b$ and $s$. CER and OCR confidence come from Tesseract.}
  \label{tab:raw}
  \begin{tabular}{@{}lccc@{}}
    \toprule
                                   & Neat   & Average & Messy  \\
    \midrule
    Baseline angle SD $\sigma_\theta$ & 0.41$^\circ$ & 0.41$^\circ$ & 0.89$^\circ$ \\
    Band spacing SD (px)           & 26.3   & 24.8    & 23.6   \\
    Baseline consistency $b$       & 0.9729 & 0.9728  & 0.9406 \\
    Spatial organization $s$       & 0.9972 & 0.9975  & 0.9977 \\
    Text fragmentation $f$         & 0.8397 & 0.8350  & 0.8599 \\
    Combined score                 & 0.9245 & 0.9226  & 0.9197 \\
    \midrule
    CER                            & 67.7\% & 75.6\%  & 77.2\%  \\
    OCR confidence (0--100)        & 54.6   & 79.9    & 63.9    \\
    Ground-truth characters        & 997    & 463     & 845     \\
    \bottomrule
  \end{tabular}
\end{table}

Baseline consistency had the widest spread of the three metrics (0.941 to
0.973) and was the only metric that placed the messy page last. Text
fragmentation placed the messy page first (least fragmented), and spatial
organization varied by less than 0.001 across the pages.

% ===========================================================================
% VI. DISCUSSION  (comments 41-47; checklist item 7)
% ===========================================================================

\section{Discussion}

With this research, we have started to address our goal of assessing the
machine readability of handwritten lab notebooks before transcription. The
readability score, computed without the OCR tool, ranked the three pages in
the same order as their CER and as the visual assessment (Table~\ref{tab:raw}).

While the score ranked the pages in the right order, every page was rated High
readability despite CERs of 68\% to 77\%. This was our biggest surprise, as
identifying scans that machines may be unable to read is the main use case.
Two of the three metrics also behaved unexpectedly. Spatial organization
barely changed across the pages (0.9972 to 0.9977). Its reference value,
a spacing standard deviation of 495 pixels, is about twenty times larger than
the spacing variation actually measured (24 to 26 pixels), so any plausible
page scores close to 1. In addition, all three notebooks are written on
printed grid paper, and the projection profile appears to detect the grid
rows along with the handwriting (on the average page it found 64 bands spaced
about 30 pixels apart), which makes line spacing look similar on every page.
A reference tied to the page's own line spacing, such as the ratio of the
spacing standard deviation to the mean spacing, together with removing the
printed grid before measurement, should make $s$ more informative.

Text fragmentation rated the messy page as the least fragmented, which
contradicts its appearance. This is not caused by the direction of the
formula: each term increases with fragmentation and the weighted sum is
inverted. On the average page, the small-component ratio contributed 88\% of
the fragmentation index (0.145 of 0.165), while component density (0.008) and
nearest-neighbor dispersion (0.012) contributed little, because their
reference values are again much larger than the measured values. The small
components appear to be dominated by fragments of the printed grid, which
thresholding and opening break into short dashes. The messy notebook's darker
grid and hand-drawn table survive thresholding as long connected lines, which
would produce fewer small components and therefore a higher $f$.
[CONFIRM: run \texttt{src/metric\_components.py} on the neat and messy pages and
replace ``appear to'' and ``would'' with the measured values.] Removing
printed rules before measurement and recalibrating the reference values are
therefore the first changes planned for $f$. Baseline consistency ($b$) was
the metric most aligned with the visual assessment, ranking the messy page
lowest.

OCR confidence did not follow transcription accuracy on these pages. The
neat page had the lowest confidence (54.6) but the lowest error (67.7\%),
and the average page had the highest confidence (79.9) but the second-highest
error (75.6\%). With three pages this is an observation rather than a
finding, but it suggests that OCR confidence alone may not reflect
transcription accuracy and should be used with care as a readability signal.

This study has four limitations: small sample size, judgement-based weights,
uncalibrated reference values, and the use of one OCR tool. First, we tested the method on three pages, one from each notebook, so
the results cannot be generalized and this study should be seen as a first
test of the approach. Second, the weights were set by human judgement. The page ranking is sensitive to these choices: with the
weight on $s$ held at 0.20, the ranking matches CER only when the weight on
$b$ is between 0.35 and 0.80, and $f$ on its own ranks the messy page as most
readable. Third, the reference values compress all three pages into a narrow
band near the top of the scale; calibrating them on a larger set of pages
(for example, setting each reference at a high percentile of the observed
values) is needed before the High, Moderate, and Low labels are meaningful.
Lastly, Tesseract produced high error rates on every page, including the page
judged neat both visually and by the score, so future work should compare
several transcription tools.

% Section numbering: Results V, Discussion VI, Conclusion VII
% (the draft has both Discussion and Conclusion as "V").
